% GENERATED FILE. Edit canonical lessons, then run npm run generate:books.
% canonical-source-hash: fnv1a64:0f09efc840931f8a
% canonical-lessons: MR-C195-savadh-karne, MR-C195-purne, MR-C195-jalne, MR-C195-thothavne, MR-C195-tulna-karne

\chapter{To warn, To be enough, To burn, To knock, To compare}
\label{ch:mr-to-warn-to-be-enough-to-burn-to-knock-to-compare}
\hlchaptermodality{195}

\begin{chapteropening}
\textbf{By the end of this chapter:} \emph{I can say to warn, to be enough, to burn, to knock and to compare.}

Complete the last lesson of chapter 195: I can say to warn, to be enough, to burn, to knock and to compare.
\end{chapteropening}

\section[\texorpdfstring{\mr{सावध} \mr{करणे} (sāvadh karṇe)}{सावध करणे (sāvadh karṇe)}]{\mr{सावध} \mr{करणे} (sāvadh karṇe) — to warn}

\label{lesson:MR-C195-savadh-karne}



\begin{quote}
Before the new one: say the Marathi for to protect, then the Marathi for to die.
\end{quote}

\subsection*{You'll want to know: \mr{सावध} \mr{करणे}}
\textbf{\mr{सावध} \mr{करणे}} — \emph{sāvadh karṇe} — \textquotedblleft{}to warn\textquotedblright{}.

\begin{grammarlens}[title={What you've built}]
One more word for this chapter.
\end{grammarlens}

\subsection*{Your turn}
\begin{itemize}
\raggedright
  \item \emph{Say it:} \emph{sāvadh karṇe}
  \item \emph{Say it:} \emph{sāvadh karṇe}, once more
  \item \emph{Say it:} say \emph{sāvadh karṇe}
  \item \emph{Recall:} say the Marathi for to protect, then the Marathi for to die, then say \emph{sāvadh karṇe} again
\end{itemize}

\begin{tcolorbox}[breakable,colback=teal!4,colframe=teal!35!black,title={Before you move on}]
What does \mr{सावध} \mr{करणे} mean? (\textquotedblleft{}To warn\textquotedblright{}.) Say it once more.
\end{tcolorbox}
\section[\texorpdfstring{\mr{पुरणे} (purṇe)}{पुरणे (purṇe)}]{\mr{पुरणे} (purṇe) — to be enough}

\label{lesson:MR-C195-purne}



\begin{quote}
Before the new one: say the Marathi for to die, then the Marathi for to warn.
\end{quote}

\subsection*{You'll want to know: \mr{पुरणे}}
\textbf{\mr{पुरणे}} — \emph{purṇe} — \textquotedblleft{}to be enough\textquotedblright{}.

\begin{grammarlens}[title={What you've built}]
One more word for this chapter.
\end{grammarlens}

\subsection*{Your turn}
\begin{itemize}
\raggedright
  \item \emph{Say it:} \emph{purṇe}
  \item \emph{Say it:} \emph{purṇe}, once more
  \item \emph{Say it:} say \emph{purṇe}
  \item \emph{Recall:} say the Marathi for to die, then the Marathi for to warn, then say \emph{purṇe} again
\end{itemize}

\begin{tcolorbox}[breakable,colback=teal!4,colframe=teal!35!black,title={Before you move on}]
What does \mr{पुरणे} mean? (\textquotedblleft{}To be enough\textquotedblright{}.) Say it once more.
\end{tcolorbox}
\section[\texorpdfstring{\mr{जळणे} (jaḷṇe)}{जळणे (jaḷṇe)}]{\mr{जळणे} (jaḷṇe) — to burn}

\label{lesson:MR-C195-jalne}



\begin{quote}
Before the new one: say the Marathi for to warn, then the Marathi for to be enough.
\end{quote}

\subsection*{You'll want to know: \mr{जळणे}}
\textbf{\mr{जळणे}} — \emph{jaḷṇe} — \textquotedblleft{}to burn\textquotedblright{}.

\begin{grammarlens}[title={What you've built}]
One more word for this chapter.
\end{grammarlens}

\subsection*{Your turn}
\begin{itemize}
\raggedright
  \item \emph{Say it:} \emph{jaḷṇe}
  \item \emph{Say it:} \emph{jaḷṇe}, once more
  \item \emph{Say it:} say \emph{jaḷṇe}
  \item \emph{Recall:} say the Marathi for to warn, then the Marathi for to be enough, then say \emph{jaḷṇe} again
\end{itemize}

\begin{tcolorbox}[breakable,colback=teal!4,colframe=teal!35!black,title={Before you move on}]
What does \mr{जळणे} mean? (\textquotedblleft{}To burn\textquotedblright{}.) Say it once more.
\end{tcolorbox}
\section[\texorpdfstring{\mr{ठोठावणे} (ṭhoṭhāvṇe)}{ठोठावणे (ṭhoṭhāvṇe)}]{\mr{ठोठावणे} (ṭhoṭhāvṇe) — to knock}

\label{lesson:MR-C195-thothavne}



\begin{quote}
Before the new one: say the Marathi for to be enough, then the Marathi for to burn.
\end{quote}

\subsection*{You'll want to know: \mr{ठोठावणे}}
\textbf{\mr{ठोठावणे}} — \emph{ṭhoṭhāvṇe} — \textquotedblleft{}to knock\textquotedblright{}.

\begin{grammarlens}[title={What you've built}]
One more word for this chapter.
\end{grammarlens}

\subsection*{Your turn}
\begin{itemize}
\raggedright
  \item \emph{Say it:} \emph{ṭhoṭhāvṇe}
  \item \emph{Say it:} \emph{ṭhoṭhāvṇe}, once more
  \item \emph{Say it:} say \emph{ṭhoṭhāvṇe}
  \item \emph{Recall:} say the Marathi for to be enough, then the Marathi for to burn, then say \emph{ṭhoṭhāvṇe} again
\end{itemize}

\begin{tcolorbox}[breakable,colback=teal!4,colframe=teal!35!black,title={Before you move on}]
What does \mr{ठोठावणे} mean? (\textquotedblleft{}To knock\textquotedblright{}.) Say it once more.
\end{tcolorbox}
\section[\texorpdfstring{\mr{तुलना} \mr{करणे} (tulnā karṇe)}{तुलना करणे (tulnā karṇe)}]{\mr{तुलना} \mr{करणे} (tulnā karṇe) — to compare}

\label{lesson:MR-C195-tulna-karne}



\begin{quote}
Before the new one: say the Marathi for to burn, then the Marathi for to knock.
\end{quote}

\subsection*{You'll want to know: \mr{तुलना} \mr{करणे}}
\textbf{\mr{तुलना} \mr{करणे}} — \emph{tulnā karṇe} — \textquotedblleft{}to compare\textquotedblright{}.

\begin{grammarlens}[title={What you've built}]
One more word for this chapter.
\end{grammarlens}

\subsection*{Your turn}
\begin{itemize}
\raggedright
  \item \emph{Say it:} \emph{tulnā karṇe}
  \item \emph{Say it:} \emph{tulnā karṇe}, once more
  \item \emph{Say it:} say \emph{tulnā karṇe}
  \item \emph{Recall:} say the Marathi for to burn, then the Marathi for to knock, then say \emph{tulnā karṇe} again
\end{itemize}

\begin{tcolorbox}[breakable,colback=teal!4,colframe=teal!35!black,title={Before you move on}]
What does \mr{तुलना} \mr{करणे} mean? (\textquotedblleft{}To compare\textquotedblright{}.) Say it once more.
\end{tcolorbox}
